\documentclass[instructions.tex]{subfiles}
\begin{document}
 
\chapter{De la vertu de pénitence.}
 
\primo Il n’y a rien au monde qui console plus pendant la vie, et qui rassure plus à la mort que la pénitence. C’est par elle que nous expions le péché, que nous nous en préservons, et que nous persévérons dans la grâce; sans elle point de salut à espérer\footnote{Nisi pœnitentiam egeritis, omnes similiter peribitis. Luc. 13. 3.}.
 
Voulez-vous, dit saint Augustin, n’être pas puni de Dieu? punissez-vous vous-même\footnote{Vis ut non puniat te Deus? te puni. S. Aug.}. Il faut que le péché soit puni par celui qui l’a commis, ou par celui contre qui il a été commis: en ce monde, par la pénitence; ou en l’autre, par les flammes. Vous êtes homme pour travailler; et vous êtes chrétien pour être pénitent\footnote{Christiana vita perpetua pœnitentia. Conc. Trid.}.
 
Vous ne serez véritablement chrétien, et vous ne serez prédestiné, qu’autant que vous serez conforme à Jésus-Christ. Sa vie a été une pénitence et un martyre continuel: voilà le chemin qu’il a tracé pour arriver à la gloire, les saints l’ont suivi. Les uns, dit saint Paul, ont souffert les moqueries, les fouets, les prisons; les autres ont été sciés, lapidés, ont passé par le fer et le feu, ont été mis à mort; d’autres ont passé leur vie en gémissant. Voilà ce qu’ont fait, pour expier leurs péchés, des âmes justes que le monde n’était pas digne de posséder.
 
Si le Fils de Dieu, si l’innocent a répandu des larmes et du sang pour des péchés qu’il n’avait point commis; si tant d’ames justes ont crucifié leur chair; que devriez-vous faire, vous qui, bien loin d’être juste, êtes chargé de tant de crimes? Refuser de faire pénitence, est une conduite plus injurieuse à Dieu, dit saint Chrysostome, que le péché même. L’impénitence est le seul crime que Dieu ne pardonne point, dit saint Jérôme\footnote{Solum crimen est quod veniam consequi non potest. Ad Sab.}.
 
\secundo Que faut-il faire pour être pénitent? Il faut imiter un vigilant économe, qui se sert de tout pour apaiser ses créanciers et payer ses dettes. Oh! que de dettes nous avons contractées envers Dieu!
 
On commence à satisfaire à Dieu, et on l’apaise, lorsqu’on se repent de ses péchés, qu’on cesse de l’offenser, qu’on change de vie. Pleurer les péchés qu’on a commis; n’en plus commettre qui méritent d’être pleurés, voilà, dit saint Grégoire, une pénitence assurée. Sans ce changement de vie, la pénitence est vaine\footnote{Ubi emendatio nulla, vana pœnitentia. S. Greg.}.
 
De quoi sert de frapper votre poitrine, de pleurer vos péchés, si vous ne voulez pas les quitter? Ce n’est pas là apaiser Dieu, ce n’est pas décharger votre conscience, mais l’endurcir et la paver de crimes, dit saint Augustin; c’est imiter ces ouvriers qui enfoncent à grands coups le pavé, pour le rendre plus dur\footnote{Conscientiam pavimentare, S. Aug.}. Vous faites de grandes aumônes: vous y êtes obligé, si vous avez de grands biens; mais l’aumône seule ne suffit pas pour apaiser le ciel. Dieu ne se laisse pas gagner par l’argent, dit Salvien, mais par la réformation des mœurs.
 
Vous vous abusez donc, si vous croyez être pénitent sans changement de vie; c’est à quoi il faut vous résoudre. Car, après tout, quel repos pouvez-vous avoir en vivant toujours ennemi de Dieu? Ce n’est pas même encore assez, pour être pénitent, de quitter le péché; il faut vivre dans le regret, dans l’humilité et la mortification, comme nous dirons ci-après.
 
\section{Résolutions}
 
\primo Quelle pénitence ai-je faite jusqu’ici? Il est bien temps que je commence.
 
\secundo Je vais donc, dès ce moment, travailler à changer de vie, faisant tous mes efforts pour ne plus retomber dans le péché. 

\prayer{C’est vous, ô Dieu! qui changez les cœurs par votre grâce; je vous présente le mien, afin que vous en arrachiez l’affection au péché, et que vous l’enflammiez d’amour pour la vertu.}
 
\end{document}
